\chapter{Additional Figures and Supplementary Artifacts}

This appendix summarizes additional figures and evaluation artifacts present in the repository that support the analyses reported in the main text. Some of these figures are already included in the body of the thesis, while others are better retained as supplementary material in order to keep the central narrative focused. The purpose of this appendix is therefore twofold: first, to document what additional visual and tabular artifacts exist in the project; second, to make explicit which analyses are directly supported by committed repository outputs rather than by ad hoc post hoc reconstruction.

\section{Discrimination Figures}

The principal all-model discrimination figure is the ROC comparison plot:

\begin{itemize}
    \item \texttt{artifacts/evaluation/roc\_curve\_all\_models.png}
\end{itemize}

This figure shows the receiver operating characteristic curves for the four headline models on the temporal test set. It is suitable for the main text because it provides a compact visual comparison of the full ranking behavior of all model families. A corresponding all-model precision--recall plot was not found as a committed artifact in the repository.

\section{Calibration Figures}

Calibration is supported by both combined and per-model reliability plots. The combined figure is:

\begin{itemize}
    \item \texttt{artifacts/evaluation/calibration\_stronger\_lr\_v3/reliability\_diagram\_all\_models.png}
\end{itemize}

Per-model reliability figures are also available:

\begin{itemize}
    \item \texttt{clinical\_logistic\_strong\_v2\_reliability.png}
    \item \texttt{clinical\_xgb\_strong\_v2\_reliability.png}
    \item \texttt{image\_stronger\_lr\_v3\_reliability.png}
    \item \texttt{multimodal\_stronger\_lr\_v3\_reliability.png}
\end{itemize}

The combined reliability diagram is most appropriate for the main text because it enables direct visual comparison. The single-model calibration plots are more naturally placed in an appendix, where they can support readers who want more detailed model-specific calibration views.

Supporting files also include:
\begin{itemize}
    \item \texttt{calibration\_metrics.json}
    \item \texttt{calibration\_summary.csv}
    \item per-model bin tables (\texttt{*\_bins.csv})
\end{itemize}

\section{Decision Curve Analysis}

Decision curve analysis is represented by two main figures:

\begin{itemize}
    \item \texttt{artifacts/evaluation/dca/decision\_curve.png}
    \item \texttt{artifacts/evaluation/dca/decision\_curve\_standardized.png}
\end{itemize}

The standard decision curve plot reports net benefit as a function of threshold probability, while the standardized variant rescales this quantity for easier comparison across settings. In the thesis, one of these figures may be placed in the main text, while the other is well suited for the appendix.

Additional supporting DCA files include:
\begin{itemize}
    \item \texttt{summary.json}
    \item \texttt{decision\_curve\_all\_models.csv}
    \item per-model threshold curves
    \item \texttt{treat\_all\_curve.csv}
    \item \texttt{treat\_none\_curve.csv}
\end{itemize}

These files are useful for reproducing figure values or extending the text-level DCA interpretation. They also make it possible to verify that the threshold-dependent utility claims in the main text are derived from script-generated evaluation outputs rather than manual calculations.

\section{Prediction-Distribution Figures}

For each headline model, the repository contains a prediction histogram showing the class-conditional distribution of predicted probabilities on the test set:

\begin{itemize}
    \item \texttt{artifacts/evaluation/prediction\_behavior\_clinical\_logistic\_strong\_v2/prediction\_histogram.png}
    \item \texttt{artifacts/evaluation/prediction\_behavior\_clinical\_xgb\_strong\_v2/prediction\_histogram.png}
    \item \texttt{artifacts/evaluation/prediction\_behavior\_image\_stronger\_lr\_v3/prediction\_histogram.png}
    \item \texttt{artifacts/evaluation/prediction\_behavior\_multimodal\_stronger\_lr\_v3/prediction\_histogram.png}
\end{itemize}

These figures are not essential to the main narrative, but they are useful supplementary diagnostics because they show how strongly each model separates positive and negative examples at the score-distribution level. They are therefore more appropriate for an appendix than for the main results chapter.

Each prediction-behavior directory also contains supporting tabular files such as:
\begin{itemize}
    \item \texttt{summary.csv}
    \item \texttt{prediction\_distribution\_by\_class.csv}
    \item \texttt{top\_false\_positives.csv}
    \item \texttt{top\_false\_negatives.csv}
    \item \texttt{predictions\_copy.csv}
\end{itemize}

\section{Qualitative Grad-CAM Assets}

The repository contains structured Grad-CAM outputs for validation-set true positives, false positives, and false negatives of the image-only model. These are stored in:

\begin{itemize}
    \item \texttt{artifacts/interpretability/gradcam\_val\_tp/}
    \item \texttt{artifacts/interpretability/gradcam\_val\_fp/}
    \item \texttt{artifacts/interpretability/gradcam\_val\_fn/}
\end{itemize}

Each stored PNG is a three-panel figure containing the original chest radiograph, a heatmap, and an overlay image. According to the repository metadata, these examples are selected from validation predictions and organized by error type. The available files include 10 true positives, 20 false positives, and 10 false negatives.

A practical subset for appendix use consists of:
\begin{itemize}
    \item \texttt{tp\_01\_study\_51021295.png}
    \item \texttt{tp\_02\_study\_53631275.png}
    \item \texttt{fp\_01\_study\_57560288.png}
    \item \texttt{fn\_01\_study\_54411196.png}
    \item \texttt{fp\_02\_study\_52432239.png}
\end{itemize}

These cases were recommended from repository ordering itself rather than from external clinical ranking. They are therefore suitable as representative supplementary examples rather than as statistically privileged cases. More broadly, the Grad-CAM outputs should be interpreted as illustrative qualitative diagnostics. They are selected by score and error mode on held-out predictions and are not inferential estimates of causal model behavior.

\section{Model Histories and Missing Supplementary Figures}

Training-history PNG files were not found as committed artifacts. However, several neural run directories contain \texttt{history.json} and \texttt{summary.json} files, including:

\begin{itemize}
    \item \texttt{artifacts/models/image\_multilabel\_pretrain\_densenet121\_strong\_v2/history.json}
    \item \texttt{artifacts/models/image\_pneumonia\_finetune\_densenet121\_u\_ignore\_temporal\_stronger\_lr\_v3/history.json}
    \item \texttt{artifacts/models/multimodal\_pneumonia\_densenet121\_triage\_u\_ignore\_temporal\_stronger\_lr\_v3/history.json}
\end{itemize}

These could be plotted in future if additional training-curve figures are desired, but they are not currently stored as ready-to-include image assets. Likewise, final \texttt{u\_zero} evaluation figures for the headline runs were not found in the evaluation directories, even though \texttt{u\_zero} manifests exist.

\section{Script-to-Artifact Provenance Mapping}

The supplementary artifacts listed above are generated through a dedicated evaluation and interpretability layer. The main script-to-artifact relationships are summarized in Table~\ref{tab:appendix_artifact_provenance}.

\renewcommand{\arraystretch}{1.15}
\small
\begin{longtable}{p{5.0cm} p{8.5cm}}
\caption{Compact mapping between repository scripts and generated supplementary artifacts.}
\label{tab:appendix_artifact_provenance} \\
\toprule
\textbf{Script / Module} & \textbf{Typical Outputs} \\
\midrule
\endfirsthead

\toprule
\textbf{Script / Module} & \textbf{Typical Outputs} \\
\midrule
\endhead

\texttt{src/evaluation/bootstrap\_eval.py} &
Bootstrap JSON summaries for AUROC/AUPRC confidence intervals and paired model deltas. \\

\texttt{src/evaluation/calibration\_analysis.py} &
\texttt{reliability\_diagram\_all\_models.png}, per-model reliability plots, \texttt{calibration\_metrics.json}, \texttt{calibration\_summary.csv}, bin tables. \\

\texttt{src/evaluation/decision\_curve\_analysis.py} &
\texttt{decision\_curve.png}, \texttt{decision\_curve\_standardized.png}, \texttt{decision\_curve\_all\_models.csv}, per-model threshold curves, \texttt{summary.json}. \\

\texttt{scripts/check\_prediction\_behavior.py} &
Per-model prediction histograms, score-distribution summaries, top false-positive and false-negative tables. \\

\texttt{scripts/evaluate\_normal\_vs\_abnormal\_negatives.py} &
Hard-negative JSON summaries comparing pneumonia-versus-normal and pneumonia-versus-abnormal-negative performance. \\

\texttt{scripts/generate\_gradcam\_examples.py} &
Grad-CAM PNG panels, selection summaries, and auxiliary heatmap arrays for chosen validation-set true positives, false positives, and false negatives. \\
\bottomrule
\end{longtable}
\normalsize

This provenance mapping is important because it shows that the supplementary figures discussed in the thesis are directly tied to committed scripts and saved outputs rather than to unsupported manual figure generation.

\section{Recommended Use}

Based on the repository inventory, the most useful supplementary figures for the thesis appendix are:

\begin{itemize}
    \item the standardized decision curve plot if the standard DCA plot is already used in the main text,
    \item the per-model reliability plots,
    \item the per-model prediction histograms,
    \item an expanded qualitative Grad-CAM gallery using a few additional representative examples,
    \item calibration and DCA summary JSON/CSV files as traceable computational support.
\end{itemize}

By contrast, training-curve figures and final \texttt{u\_zero} result figures would require additional plotting or new outputs and are therefore not treated as current appendix-ready assets in this thesis. No unsupported figures are included in this appendix; all referenced artifacts correspond to files or outputs identified directly in the repository.